\begin{abstract}
红外干涉测厚通过反射光谱中的条纹反演外延层几何厚度，折射率色散、界面反射相移及高阶往返传播共同影响条纹的相位演化与线形。针对碳化硅和硅各两个入射角的实测光谱，建立由两束相位模型到多光束振幅模型的分层反演方法，区分反射率拟合误差、条件重抽样波动与材料参数引起的系统敏感性。四份附件共29876组观测，保留原始数值，对共同零值端点和262个超过100\%的反射率点分别核对，不采用插补或裁剪改变条纹结构。

针对问题一，从同一出射波前上的光程差出发，建立包含Snell折射关系、界面反射相移和波长色散的两束干涉模型。引入色散相位坐标，推导同类条纹级次差与厚度的关系，说明周期法应考虑群光学厚度；进一步分析$10^\circ$与$15^\circ$两角对折射率和厚度分离的病态性。

针对问题二，以傅里叶频率和峰位级次确定初值，在文献色散适用的$2000$--$3300\cmiv$窗口进行双角联合拟合。通过变量分离处理低阶背景与条纹包络，碳化硅共同厚度为$7.38431\um$，两角独立结果为$7.40781\um$和$7.36078\um$，反射率RMSE为0.02440个百分点。有效窗口变化给出$7.37384$--$7.39770\um$，支持报告条件厚度约$7.38\um$。

针对问题三，由多次往返振幅的几何级数得到多光束模型，区分形成条件、可见条件及数据可辨识性。硅采用固定外延层色散与有效衬底响应，共同厚度为$3.40521\um$；多光束模型将同参数两束基准的RMSE从4.36318降至0.67894个百分点，并在36个留出块中的34块降低误差。碳化硅高阶修正仅为$-0.000158\um$，分块预测误差反增3.25\%，故保留两束结果。

灵敏度分析和各120次移动块残差重抽样表明，密集采样形成的窄条件区间不能覆盖折射率及双角不一致带来的影响。最终建议在声明材料模型的前提下报告碳化硅约$7.38\um$、硅约$3.41\um$，并同步保留分角结果与参数敏感性。
\keywords{红外干涉测厚；折射率色散；多光束干涉；变量分离；分块验证；条件不确定性}
\end{abstract}
